%% one_way_relu_part2.tex — Part II: the research program the Part I
%% results opened. Status: prospectus with pre-registered designs;
%% instruments under construction.
\documentclass[11pt]{article}
\usepackage{ruachtov}
\renewcommand{\rtproject}{Backpropagation}
\rtdate{17 November 2026}
\usepackage[margin=1in]{geometry}
\usepackage{amsmath}
\usepackage{amssymb}
\usepackage{graphicx}
\usepackage{float}
\usepackage{booktabs}
\usepackage{caption}
\captionsetup{labelfont=bf}

\newcommand{\relu}{\mathrm{ReLU}}
\newcommand{\modelEmnistNet}{\textsc{EmnistNet}\xspace}
\newcommand{\modelCifarNet}{\textsc{CifarNet}\xspace}
\newcommand{\modelVggEleven}{\textsc{Vgg-11}\xspace}

\author{Heath\thanks{With Bocher and other AI Agents of Ruach Tov
Collective (\texttt{https://ruachtov.ai}).}}

\begin{document}
\rttitle{One-way ReLU, Part II: Unit Telemetry and the
Decomposition of LayerNorm}{A research program --- pre-registered
designs for the questions Part~I opened}

\begin{abstract}
Part I~\cite{partone} established that the backward convention at the rectifier's
kink is a free experimental parameter with large, regime-dependent
effects: one-way training is more accurate where units squelch
(\modelEmnistNet{} $+0.40$, CIFAR-10 $+2.61$), less accurate where
squelching serves as implicit capacity pruning (\modelVggEleven{}
$-3.1$ under Adam), and interacts with LayerNorm so strongly that the
normalizer stably reverses the direction of its effect on a small
transformer. Part II opens with that transformer result in full --- withheld
from Part I pending the decomposition it demands --- and lays out
the program the Part I results open. We propose:
(1) \emph{unit telemetry} --- instrumented kernels that detect
the moment a unit would go silent and record the circumstances,
enabling a causal trial of whether the silenced condition
\emph{spreads} through a network, the directions of its spreading
mechanisms, and whether single-layer one-way reconditioning
\emph{interrupts} the spread; (2) a direct study of when
silenced units indicate a problem and when they are useful pruning,
including the dropout interaction; and (3) a decomposition of
LayerNorm into three primitives --- mean-shift, symmetric
deviation-scale, and a novel \emph{balanced} (per-side) deviation
scale built on the positive/negative semideviation --- factorially
recombined to locate which ingredient reverses one-way's effect.
\end{abstract}

\section{First result: removing LayerNorm and adding one-way on WikiText-2}
\label{sec:wikitext}

Part I moved from characters (EMNIST) to natural images
(CIFAR-10); this experiment moves off images and off convolutions
entirely: a small transformer language model (two layers,
$d_{\mathrm{model}}{=}128$, four heads, FFN width 512, tied
embeddings) on word-level WikiText-2, validation perplexity as the
metric. The transformer's feed-forward block ($\max(0, xW_1)W_2$) is
the single most-executed ReLU site in modern computing. Table~\ref{tab:wikitext} gives the full factorial at $n{=}10$
paired seeds per cell:

\begin{table}[H]
\centering
\caption{WikiText-2 validation perplexity (lower is better),
$n{=}10$ paired seeds per cell, Adam throughout. LN$^{+}$/LN$^{-}$
denotes LayerNorm present/removed (all LayerNorms replaced by
identity).}
\label{tab:wikitext}
\begin{tabular}{lcc}
\toprule
 & LN$^{+}$ & LN$^{-}$ \\
\midrule
ReLU     & $387.66 \pm 8.40$ & $331.17 \pm 7.58$ \\
one-way  & $397.54 \pm 6.89$ & $\mathbf{327.02 \pm 6.68}$ \\
\bottomrule
\end{tabular}
\end{table}

\textbf{Replacing the standard recipe (LN$^{+}$ReLU) with the
renegade recipe (LN$^{-}$one-way) improves perplexity by $60.6$
points --- $15.6\%$ --- $t{=}{-}18.1$, on all 10 of 10 paired
seeds.} The decomposition matters more than the total, and the two
factors contribute very unequally. \emph{Nearly all of the
improvement comes from removing LayerNorm}: that step alone makes
plain ReLU better by $56.5$ points ($t{=}{-}16.9$, 10/10). The
literature offers no direct precedent --- LayerNorm predates the
Transformer and was baked in from the architecture's birth, so the
train-from-scratch ablation at small scale appears never to have
been reported (recent work removes LN only at inference time, from
models trained with it). One-way's own increment on the LN$^{-}$
substrate is comparatively small --- $4.2$ points, about $1\%$,
roughly half a seed-level standard deviation ($t{=}{-}3.2$, 8/10
seeds) --- large enough to establish the \emph{sign} of the effect
at this $n$, not large enough to claim its \emph{magnitude} with
confidence; a higher-powered run is in progress. What the factorial
does establish firmly is the interaction: under LN$^{+}$, one-way
makes the model \emph{worse} by $9.9$ points ($t{=}15.0$, 0/10) ---
while under LN$^{-}$ it makes the model better. LayerNorm stably
reverses the direction of one-way's effect. Notably, no FFN channel ever squelches in \emph{any} cell
(silenced fraction $0.000$ throughout, with or without LN --- the
residual stream alone suffices at this depth), so both effects here
are a \emph{second mechanism}, distinct from the squelch-repair
mechanism of Part I: in the un-normalized regime, blocking
suppression-direction traffic evidently reduces gradient
interference; under re-centering, the same blocked traffic was
carrying useful mean-shift information.

Scope, stated plainly: two layers, $d{=}128$, one dataset ---
the small-scale regime, where LN's depth-emergent stabilization has
nothing to stabilize. The depth at which LN begins to help
--- and whether one-way moves that frontier --- is the open question
this result makes precise --- and the decomposition of LayerNorm
into mean-shift and deviation-scale primitives (including a
balanced, per-side scale built on the semideviation), factorially
recombined against both conventions, is the opening experiment of
Part~II (a momentum-SGD arm at this scale trains
nothing at all in either convention: the un-normalized transformer
is Adam-dependent, a reconditioning interaction consistent with
Section~3.3's findings).

\section{A nomenclature for component genotypes}
\label{sec:nomenclature}

Experimental deep learning has no standard notation for the thing it
does most: training otherwise-identical models that differ in one
named component. The literature writes ``w/o LN,'' ``ablated,''
``-BN,'' ``ours,'' and full prose table rows --- every paper its own
convention, none composable. Genetics solved this problem a century
ago, and we propose borrowing its solution directly.

\subsection{The borrowed grammar}

In genetic nomenclature, a superscript minus marks a
\emph{loss-of-function} allele ($a^{-}$), a superscript plus marks
\emph{wild type} ($a^{+}$), and a knockout organism is written
$\mathit{Gene}^{-/-}$. The genotype string
$\mathit{Kit}^{W\text{-}v}$ names an \emph{allele} --- a specific
variant at a named locus --- in its superscript. We map this onto
compute graphs:

\begin{itemize}
\item \textbf{Knockout:} $X^{-}$ --- component $X$ removed from the
graph, everything else identical. Our LN$^{-}$ transformer: every
LayerNorm replaced by identity. (Genetics: $a^{-}$.)
\item \textbf{Wild type:} $X^{+}$ --- the component as the standard
recipe ships it. LN$^{+}$, Adam$^{+}$. Marked explicitly rather than
left implicit, because the factorial's control cell is a measurement,
not a default. (Genetics: $a^{+}$.)
\item \textbf{Engineered allele:} $X^{\mathit{name}}$ --- a
\emph{designed variant at the same site}: neither absence nor stock,
but modified function. ReLU$^{\mathit{ow}}$ is the one-way backward
convention at a ReLU site; Adam$^{\mathit{sg}}$ (sign-gated
accumulator) and Adam$^{\mathit{sr}}$ (sign-reset on revival) are
alleles of Adam; LN$^{\mathit{eq}}$ is LayerNorm with
equinormalized scale (Section~\ref{sec:ln-decomposition}).
(Genetics has no symmetric standard here --- gain-of-function and
engineered insertions are handled by ad-hoc named superscripts and
transgene notation; the field is arguably on the cusp of needing
this slot, while machine learning, which designs new function
routinely, needs it today.)
\item \textbf{Genotype strings:} components concatenate.
LN$^{-}$Adam$^{+}$ReLU$^{\mathit{ow}}$ names one cell of a factorial
completely; a $2\times3$ design is the cross
LN$^{\{+,-\}} \times$ Adam$^{\{+,\mathit{sg},\mathit{sr}\}}$.
\item \textbf{Site scoping:} where a component appears at many
sites, a site list subscripts the allele:
ReLU$^{\mathit{ow}}_{c1,c2}$ is one-way applied at the first two
convolutional sites only (1-indexed from the input; the per-site
placement arms of Part I).
\end{itemize}

\subsection{Why the analogy is load-bearing, not decorative}

The mapping works because the experimental situations are genuinely
isomorphic. A paired-seed ablation \emph{is} a knockout study: same
genome (the declarative network specification), one locus modified,
phenotype measured against wild type. In our toolchain the analogy
is exact enough to diff: the one-way and strict arms of Part I
differ in three facts of one Prolog file --- the genotype delta is
six printed lines. The notation then buys what genetics bought:
\begin{enumerate}
\item \textbf{Compression with precision.} Table headers become
genotype strings; ``our method'' and ``baseline'' --- which hide the
intervention --- are replaced by names that state it.
\item \textbf{Composability.} Factorial designs read as crosses;
interaction claims (``LN$^{+}$ reverses ReLU$^{\mathit{ow}}$'')
name both loci.
\item \textbf{A place for engineered function.} The named-allele
slot distinguishes \emph{removing} a mechanism from \emph{replacing}
it --- the distinction between LN$^{-}$ and LN$^{\mathit{eq}}$ that
``w/o LN'' cannot express.
\end{enumerate}

We use this notation throughout the remainder of this paper and
retroactively in discussing Part I, and we offer it to the field for
evaluation: it costs one sentence to teach, inherits a century of
precedent, and composes.

\section{Unit telemetry: mapping the silenced condition as it
spreads}
\label{sec:cartography}

Part I measured squelching as a static census (what fraction of units
are silent at epoch's end). The sharper instrument is a
\emph{cartograph}: record each silencing event at the moment it
happens, with its circumstances, and reconstruct the spatiotemporal
map of the condition's spread.

\subsection{Instrumentation: breakpoints at the moment of silencing}

Because the collective emits its own kernels from declarative
specifications (Part I, Section 3.1), instrumentation can be placed
\emph{inside} the backward operator rather than bolted on in Python.
Three escalating designs:

\begin{enumerate}
\item \textbf{Counting taps.} The emitted backward kernel gains a
per-unit epoch-windowed counter: when a unit's pre-activation has
been nonpositive for every input of the window \emph{and} the
arriving gradient would push it further negative, the tap increments
a ``would-silence'' counter in a side buffer. Near-zero overhead;
gives the full silencing time-series per unit.
\item \textbf{Signalling-NaN tracers (``radioactive data'').} Seed
chosen units' dormant-side values with signalling NaN payloads and
let the hardware's NaN-propagation do the tracing: any downstream
value contaminated by a traced unit carries the mark in its payload
bits. Where Part I used signed zero as a message, this uses the
NaN payload as a \emph{dye}. Caveat from Part I's bit-level work:
torch sanitizes NaN in some reductions, so the dye survives only
elementwise paths --- which is precisely where the squelching
mechanism lives.
\item \textbf{Raising breakpoints.} A branch in the emitted kernel
\texttt{raise()}s (or traps, on CPU, via floating-point exception
enablement) the first time a watched unit crosses the silencing
predicate --- a hardware-breakpoint discipline for training dynamics.
Expensive; reserved for interrogating single events discovered by
the counting taps.
\end{enumerate}

\subsection{The standing evidence: silenced inputs concentrate
upstream of silenced units}

Part I's static probe already found the association the cartograph
will resolve in time. In the CIFAR-10 testbed under Adam, for
\emph{silenced} penultimate-layer units, $4.0\%$ of incoming weight
mass originates at silenced conv-2 channels; for \emph{surviving}
units, $1.1\%$ --- a $3.6\times$ concentration. Under leaky (no
absorbing state) the asymmetry vanishes ($0.7\%$ vs $0.7\%$). The
association is real and specific to the absorbing dynamics; what it
cannot yet tell us is direction: do starving inputs silence their
consumers, do silenced consumers release their suppliers, or does a
common cause (Adam's amplified suppression traffic) silence both?

A second early observation sharpens the stakes. Call the first
convolutional layer the network's \emph{retina} --- the only layer
that touches the raw stimulus. (Layer numbering convention,
standardized here and throughout: conv sites are numbered
c1, c2, \ldots{} \emph{from the input}, 1-indexed, matching the
\texttt{op/4} facts of Part I Listing 1; \modelVggEleven{} has
conv sites c1--c8 followed by fc sites f1--f2. The experiment
harness's internal 0-indexed site names are mapped accordingly in
the published analysis scripts.) At depth four (\modelEmnistNet{},
\modelCifarNet{}) the retina never silences: $0.0\%$ of c1
channels in every run of every arm, plausibly protected by
weight-sharing's many chances to fire. At depth eleven
(\modelVggEleven{}) the retina c1 itself also never silences ---
but the protection ends there: \emph{c2, the second convolution},
silences at $64\%$ on average --- in individual seeds $86\%$ ---
among the worst conv sites in the network, despite being second
from the stimulus. A condition that intensifies with the
number of layers \emph{above} it is at least consistent with
arrival from above --- suppression traffic amplified through ten
downstream sites accumulating on the layer least able to re-fire
once its input statistics are fixed. The retina may be where the
condition \emph{finishes}, not where it starts. The telemetry's
event times will say.

\subsection{The causal trial}
\label{sec:causal}

The cartograph converts the association into a trial with three
registered questions:

\begin{enumerate}
\item \textbf{Does it spread?} With event times for every silencing,
test whether a unit's hazard of silencing increases after the
silencing of its suppliers (upstream conditioning) beyond what the
global silencing rate predicts.
\item \textbf{Which direction?} Compare forward spread (suppliers
first, consumers follow: starvation) against backward spread
(consumers first, suppliers follow: gradient-path closure --- when a
consumer silences, its suppliers lose one gradient route and their
remaining signal weakens). The two mechanisms have opposite
temporal signatures.
\item \textbf{Does single-layer one-way block the flow?} The
intervention arm: apply one-way at exactly one layer (the per-site
mixing of Part I) and ask whether the condition's spread, measured
upstream and downstream of the treated layer, is interrupted --- an
insulation experiment. If spread exists and one-way interrupts it,
the per-site treatment becomes a targeted intervention rather than a
whole-network convention --- and the standing placement evidence
(Section~\ref{sec:placement}) already shows per-site treatment
outperforming every uniform convention at depth.
\end{enumerate}

\textbf{Scheduled pilot.} An $n{=}3$ epoch-resolution run is
scheduled on \modelVggEleven{} (Adam, ReLU): per-unit silencing
masks recorded at every site after every epoch, together with the
interlayer weight-mass matrices, giving the first time-resolved
picture of the spread --- enough to distinguish a directed cascade
(silencings of suppliers systematically \emph{precede} those of
their consumers, or the reverse) from a dynamic equilibrium (units
exchanging silence bidirectionally around a stable aggregate rate).
The mechanism may be causal, or it may be an equilibrium; the event
order will say which. Event-level (per-batch) taps in the emitted
kernels follow if the epoch-resolution pilot finds structure.

\section{Is a silenced unit always a problem?}
\label{sec:pathology}

Part I's capacity-economics finding --- one-way more accurate at
\modelCifarNet{} scale, less accurate on \modelVggEleven{} under
Adam, where plain ReLU wins \emph{while half its network is silent}
--- forces the question. The proposed decomposition:

\begin{itemize}
\item \textbf{Silence as lost function}: at tight capacity, every
silenced unit is lost capacity; one-way restores it and accuracy
rises.
\item \textbf{Silence as pruning}: at surplus capacity, silencing is
implicit architecture search; preventing it preserves units the task
cannot use.
\item \textbf{Silence as regularization}: the dropout-adjacent
hypothesis --- stochastic unit removal during training improving
generalization. Part I's first test (\modelVggEleven{}, dropout
$0.5$ at the fc sites, $n{=}5$) \emph{falsified} the strong form:
explicit dropout did not close the one-way deficit, and ReLU's
silencing persisted undiminished through dropout --- the two
mechanisms coexist rather than substitute. The weak form (dropout
styles and sites) remains open, as does the interaction's other
face: does one-way \emph{plus} dropout outperform one-way alone,
i.e.\ does one-way need an explicit regularizer where ReLU brings
an implicit one?
\end{itemize}

The optimizer column matters throughout: under momentum-SGD at
\modelVggEleven{} depth (untuned learning rate), every arm trains
poorly, yet one-way degrades most gracefully ($0.51$ vs ReLU's
$0.39$ mean) --- the opposite ordering from Adam. Any theory of
silence-as-pruning must therefore be a theory of
\emph{Adam-mediated} pruning.

\subsection{Standing evidence: placement is not fungible}
\label{sec:placement}

Two Part I measurements already show that one-way's effects are
\emph{site-} and \emph{class-}inhomogeneous, which any mechanism
proposed here must explain.

\textbf{Per-site placement (\modelVggEleven{}, Adam, $n{=}5$
paired).} Applying one-way at only the two \emph{shallow} conv
sites (c1, c2 --- the retina and its immediate consumer; harness
arm \texttt{mix:c0,c1} in 0-indexed internal names) while the rest
of the network trains strict ReLU is \emph{more} accurate than
every uniform convention:
$+1.82$ points over all-ReLU ($t{=}4.2$, 5/5) and $+4.86$ over
all-one-way ($t{=}12.5$, 5/5). The mirrored placement --- one-way at
only the two \emph{deepest} sites (f1, f2 --- harness
\texttt{mix:f0,f1}) --- is slightly
\emph{less} accurate than all-ReLU ($-0.66$, $t{=}{-}1.8$) and
clearly worse than the shallow placement ($-2.48$, $t{=}{-}5.9$,
0/5). Protecting the retina while leaving the deep stack free to
prune is, at this depth, the most accurate convention measured ---
consistent with the Section~\ref{sec:cartography} picture in which
deep silence is (Adam-mediated) pruning but retinal silence is
unrecovered capacity: the sites where silence is a problem are
exactly the sites where prevention helps.

\textbf{Per-class distribution (\modelEmnistNet{}, Part I,
$n{=}50$/$n{=}30$).} One-way's accuracy improvement is not evenly
distributed over classes, and the \emph{shape} of the distribution
follows the optimizer: under SGD the gain is a uniform tide (all 47
of 47 classes improve, no confusion pair singled out); under Adam it
concentrates in the confusable families ($f$, $0$, $9$, $I$, $g$,
$q$ lead, with a slight negative tail) --- the repair lands where
Adam's squelching carved. Placement inhomogeneity across layers and
gain inhomogeneity across classes are plausibly the same fact seen
from two sides: one-way helps exactly where the condition is a
problem, and nowhere else.

\subsection{The active ingredient: completing the
forward/backward square}
\label{sec:square}

Part I's arms occupy three corners of a $2\times2$ design ---
\{forward: strict, leaky\} $\times$ \{backward: standard,
sign-gated\} --- leaving a confound: is one-way's accuracy a
property of the \emph{gate}, or of the forward/backward
\emph{mismatch} family it belongs to? Two measurements close the
question (all CIFAR-10 \modelCifarNet{}, Adam$^{+}$, paired seeds).

\textbf{The fourth corner.} ReLU$^{\mathit{owl}}$ (leaky forward,
sign-gated backward) at $n{=}5$: accuracy $0.7109 \pm 0.0034$ ---
$+3.01$ over ReLU$^{+}$ ($t{=}8.4$, $p{=}0.0011$), $+2.44$ over
leaky ($t{=}9.4$, $p{=}0.0007$), and statistically indistinguishable
from ReLU$^{\mathit{ow}}$ ($-0.31$, $t{=}{-}1.7$, $p{=}0.16$;
Table~\ref{tab:square}):

\begin{table}[H]
\centering
\caption{The completed square (CIFAR-10, Adam$^{+}$). The gated
column is $\sim$$+3$ points over the standard column in both rows;
the forward convention contributes at most a statistically
unresolved sliver. $n{=}10$ for the three standing cells, $n{=}5$
for ReLU$^{\mathit{owl}}$.}
\label{tab:square}
\begin{tabular}{lcc}
\toprule
 & backward standard & backward sign-gated \\
\midrule
forward strict & $0.6825$ (ReLU$^{+}$) & $\mathbf{0.7128}$ (ReLU$^{\mathit{ow}}$) \\
forward leaky  & $0.6866$ (leaky)      & $\mathbf{0.7109}$ (ReLU$^{\mathit{owl}}$) \\
\bottomrule
\end{tabular}
\end{table}

\textbf{The swap control.} The complementary confound --- that the
benefit might be ``leaky-style training followed by strict
inference'' rather than the gate --- is tested by taking
leaky-trained checkpoints and \emph{evaluating} them with strict
forward. The swap \emph{loses} accuracy on every seed: $-0.40$
points on EMNIST ($n{=}5$ paired, $t{=}{-}4.6$, $p{=}0.0099$) and
$-0.70$ on CIFAR-10 ($n{=}10$ paired, $t{=}{-}7.1$,
$p{=}0.00005$) --- leaky-trained networks use their negative-side
transmissions, and the clamp deletes computation the weights depend
on. The swap lands \emph{below plain ReLU}, while one-way sits
$+3.3$ above it (CIFAR).

Together the two controls give the attribution from both directions:
\textbf{the benefit follows the sign-gated backward wherever it
goes, and no combination of forward conventions without the gate
reproduces any of it.} The gate --- blocking suppression-direction
gradient to inactive units while admitting revival --- is the active
ingredient.

\textbf{The Adam$^{-}$ stratum.} The same four corners under
momentum-SGD (lr $10^{-2}$, $\mu{=}0.9$, $n{=}5$ paired) test
whether the attribution survives the optimizer change:
ReLU$^{+}$ $0.6928 \pm 0.0057$, leaky $0.6887 \pm 0.0050$,
ReLU$^{\mathit{ow}}$ $0.7065 \pm 0.0107$, ReLU$^{\mathit{owl}}$
$0.7037 \pm 0.0033$. The gated column again leads both rows ---
ReLU$^{\mathit{ow}}{-}$ReLU$^{+}$ $= +1.37$ ($t{=}2.5$, 4/5),
ReLU$^{\mathit{owl}}{-}$ReLU$^{+}$ $= +1.09$ ($t{=}3.3$, 5/5),
gated-vs-gated again indistinguishable ($-0.28$, $t{=}{-}0.6$) ---
while plain leaky is slightly \emph{less} accurate than ReLU$^{+}$
($-0.41$, $t{=}{-}2.1$, 0/5). Two differences from the Adam$^{+}$
stratum are themselves informative: the gate's margin shrinks from
$\sim$$+3$ to $\sim$$+1.2$ points (under SGD almost no units
squelch on this testbed, so there is less for the gate to repair),
and yet the margin does \emph{not} vanish --- the gate is more
accurate even in a regime with essentially no silenced units, which
means part of its effect is not squelch-repair at all. This echoes
the second mechanism of Section~\ref{sec:wikitext} (where FFN
squelch was zero in every cell): \emph{blocking
suppression-direction traffic appears to reduce gradient
interference generally, with squelch-repair as the amplifier that
turns a one-point effect into a three-point effect where silencing
is rampant.}

\subsection{The accumulator question: alleles of Adam}
\label{sec:adamalleles}

If the gate blocks wrong-way per-sample signals, should the
optimizer's \emph{memory} still carry and apply them? Stock Adam
does: while a unit is silenced its rows' first moment continues
applying pre-silence suppression momentum, and the second moment
decays toward zero, amplifying whatever signal next passes (Part I's
finding that revival-step magnitudes are nearly
$\alpha$-independent under Adam is this amplification measured).
Two repair alleles: Adam$^{\mathit{sg}}$ (sign-gated accumulator:
rows of silenced units have both moments frozen; the passed revival
signal steps directly) and Adam$^{\mathit{sr}}$ (sign-reset:
Adam$^{\mathit{sg}}$ plus zeroing the moments at the revival
transition --- the revived unit forgets the suppression era
entirely).

At $n{=}5$ on \modelCifarNet{}, neither allele moves final accuracy
(ReLU$^{\mathit{ow}}$: Adam$^{\mathit{sg}}{-}$Adam$^{+}$ $=$
$+0.27$, $t{=}0.6$; Adam$^{\mathit{sr}}$: $-0.09$, $t{=}{-}0.3$;
ReLU$^{+}$ likewise null). The honest reading: the accumulator's
gate-dishonesty measurably distorts \emph{unit-level dynamics}, but
at this capacity the network routes around the distortion at the
accuracy level. The sharper registered probe --- whether
Adam$^{\mathit{sg}}$ restores the $\alpha$ dose-response that stock
Adam erases --- remains to run.

\section{Decomposing LayerNorm}
\label{sec:ln-decomposition}

Section~\ref{sec:wikitext} found that LayerNorm~\cite{ba2016}
stably reverses one-way's effect on a small transformer. LayerNorm
(Section~\ref{sec:ln-decomposition}'s subject) is a composite
operation; the reversal
must live in one (or an interaction) of its parts. We decompose it
into primitives and recombine factorially.

\subsection{The primitives}

For a feature vector $x \in \mathbb{R}^d$ with mean
$\mu = \tfrac1d \sum_i x_i$:

\begin{enumerate}
\item \textbf{Mean-shift} (centering): $x_i \mapsto x_i - \mu$,
with the usual trainable bias restored afterward. The cheap
hypothesis --- unlikely, on its own, to reproduce the reversal, but
the control the factorial needs.
\item \textbf{Normalization} (pooled deviation-scale, without
shift):
$x_i \mapsto \mu + \gamma\,(x_i - \mu)/\sigma$, where $\sigma$ is
the usual standard deviation and $\gamma$ trainable. Deviations are
normalized to $\mu \pm 1$ \emph{without moving the mean}: scale
normalization strictly separated from centering.
\item \textbf{Equinormalization} (balanced, per-side
deviation-scale). Define the
\emph{standard positive deviation} and \emph{standard negative
deviation}
\[
\sigma_{+} \;=\; \Bigl(\tfrac{1}{N_{+}} \textstyle\sum_{x_i > \mu}
(x_i - \mu)^2 \Bigr)^{1/2},
\qquad
\sigma_{-} \;=\; \Bigl(\tfrac{1}{N_{-}} \textstyle\sum_{x_i < \mu}
(x_i - \mu)^2 \Bigr)^{1/2},
\]
with $N_{+}$, $N_{-}$ the counts of features above and below the
mean (these are the \emph{semideviations} of the finance
literature~\cite{markowitz} --- apparently unused as a
normalization primitive in deep learning). The balanced scale maps
\[
x_i \;\mapsto\; \mu + \gamma \cdot
\begin{cases}
(x_i - \mu)/\sigma_{+}, & x_i > \mu,\\[2pt]
(x_i - \mu)/\sigma_{-}, & x_i < \mu,
\end{cases}
\]
restoring left/right \emph{balance}: after the map, the positive and
negative sides each have unit semideviation. A skewed activation
distribution --- precisely what a rectifier manufactures --- is made
symmetric about its mean without the sides contaminating each
other's scale. Note the standard LayerNorm divides both sides by the
one pooled $\sigma$, which a skewed distribution's heavy side
dominates: in this taxonomy standard LN applies an
\emph{asymmetric} effective scaling, and the balanced primitive is
the symmetric repair.
\end{enumerate}

\subsection{The factorial}

\begin{table}[H]
\centering
\caption{The normalization factorial:
$\{\text{none}, \text{mean-shift}\} \times \{\text{none},
\text{normalization}, \text{equinormalization}\}$, each cell
$\times$ \{ReLU, one-way\} on the Part I WikiText-2 testbed. The
\emph{normalization} column is the standard pooled-$\sigma$ scale
--- asymmetry-\emph{preserving}, since a skewed distribution's heavy
side dominates the pooled deviation; \emph{equinormalization}
balances the two sides by their own semideviations. Standard
LN$^{+}$ and the Part I LN$^{-}$ cells are recovered within the
factorial.}
\label{tab:lnfactorial}
\begin{tabular}{lccc}
\toprule
 & no scale & normalization & equinormalization$^{*}$ \\
\midrule
no shift   & LN$^{-}$ (Part I) & scale-only & equinorm-only \\
mean-shift & shift-only & LN$^{+}$ (Part I) & balanced LN \\
\bottomrule
\end{tabular}

{\footnotesize $^{*}$ normalization by semideviations.}
\end{table}

Twelve cells at $n{=}10$ paired seeds resolve: which primitive
carries LayerNorm's accuracy effect at this scale; which carries the
reversal of one-way; and whether balanced scaling --- the primitive
that restores the symmetry the rectifier destroys --- behaves
differently with a backward convention that is itself about the
asymmetry of the two sides. The conjecture worth registering in
advance: \emph{mean-shift is the reverser} (re-centering turns
blocked suppression-direction gradient into lost mean-information),
and \emph{balanced scaling is the cell where one-way and
normalization finally cooperate}.

\section{Planned experiments: the decomposed atoms on
pre-transformer networks}
\label{sec:cnn-atoms}

The Part I CNNs are the natural second venue for the decomposition,
for a reason the transformer cannot offer: \textbf{the
convolutional testbeds are norm-free by construction} --- their
native state is the LN$^{-}$ corner of
Table~\ref{tab:lnfactorial}. On a transformer we ask which atom of
an entrenched normalizer causes an observed reversal; on
\modelEmnistNet{} we can ask the prior question: \emph{what does
each atom do to a network that never had any of them} --- and does
any atom reproduce, on a CNN, the reversal of one-way that full
LayerNorm produces on the transformer?

\subsection{Designs}

\begin{enumerate}
\item \textbf{Atoms on \modelEmnistNet{} ($n{=}50$ power).} Each
primitive --- mean-shift, symmetric deviation-scale, balanced
deviation-scale --- inserted as a per-site
\texttt{reconditioning(tensor(f\_f1), $\cdot$)} fact at the
penultimate pre-activation (the site where Part I measured $39.5\%$
squelching under Adam), crossed with \{ReLU, one-way\} and with
\{Adam, momentum-SGD\}. The registered questions: does any atom
\emph{prevent} squelching the way the full normalizer does on the
transformer (centering holds the pre-activation distribution astride
zero --- mean-shift alone should suffice if prevention is the
mechanism); does whichever atom prevents squelching also
\emph{reverse} one-way's sign on a CNN; and does balanced scaling
--- which symmetrizes exactly the distribution the rectifier skews
--- improve either convention where symmetric scaling does not.
\item \textbf{Atoms at depth (\modelVggEleven{}).} The same
insertion at the three sites Part I's per-layer map identified as
extreme (c2, the near-input spike; c7--c8, the deep-conv plateau;
f1, the fc maximum --- 1-indexed), one site at a time --- the normalization
analogue of the per-site convention mixing. At \modelVggEleven{}
scale the question acquires the capacity-economics dimension: if
silence there is pruning rather than lost function, an atom that prevents
silence should \emph{reduce} accuracy under Adam --- the
normalizer's classical benefit and the pruning benefit pull in
opposite directions, and the atoms let us watch them separate.
\item \textbf{Batch-norm as a control.} The pre-transformer
literature's own normalizer, decomposed the same way (its shift and
scale computed over the batch axis rather than the feature axis),
run at the same sites. If the CNN results track the \emph{axis} of
normalization rather than the atom, the story is about statistics
pooling; if they track the atom regardless of axis, the story is
about the primitive itself. The conjecture, extended from
Section~\ref{sec:ln-decomposition}: the reversal follows mean-shift
wherever it goes, on any axis, in any architecture.
\item \textbf{The cartograph watches.} Every arm above runs under
the Section~\ref{sec:cartography} counting taps, so each atom's
effect arrives with its silencing time-series: not only
\emph{whether} mean-shift prevents squelching on a CNN but
\emph{when} --- whether it suppresses the early-epoch silencing
burst Part I observed ($52\%$ silenced after a single CIFAR epoch) or
merely accelerates late recovery.
\end{enumerate}

\section{Status}

Part I is complete and its artifacts are public. The component-
genotype nomenclature (Section~\ref{sec:nomenclature}) is in use
throughout this document and offered for field evaluation; the
completed forward/backward square (Section~\ref{sec:square}) and
the Adam alleles (Section~\ref{sec:adamalleles}) are its first
worked applications. The Part II
instruments --- the counting taps of the cartograph
(Section~\ref{sec:cartography}), whose causal trial
(Section~\ref{sec:causal}) and silence taxonomy
(Section~\ref{sec:pathology}) consume its event streams, and the
normalization primitives of Table~\ref{tab:lnfactorial} as
\texttt{reconditioning/2} facts and harness modules (whose
pre-transformer campaign Section~\ref{sec:cnn-atoms} registers) --- are under
construction on the same declarative toolchain. Designs in this
document are pre-registered in the Part I sense: arms, metrics, and
seed counts stated before the runs.

\begin{thebibliography}{9}
\bibitem{partone} Heath, with the Ruach Tov Collective.
\emph{One-way ReLU: a sign-dependent backward convention for the
rectifier.} Part I. Ruach Tov, November 2026.
\bibitem{markowitz} H.~Markowitz.
\emph{Portfolio Selection: Efficient Diversification of
Investments.} Wiley, 1959. (Semivariance.)
\bibitem{ba2016} J.~L.~Ba, J.~R.~Kiros, G.~E.~Hinton.
\emph{Layer Normalization.} arXiv:1607.06450, 2016.
\end{thebibliography}

\end{document}
